% Part: first-order-logic
% Chapter: sequent-calculus
% Section: quantifier-rules

\documentclass[../../../include/open-logic-section]{subfiles}

\begin{document}

\olfileid{fol}{seq}{qrl}

\olsection{قواعد المكممات}

\subsection{قواعد $\lforall$}

\begin{defish}
\Axiom$ !A(t), \Gamma \fCenter \Delta$
\RightLabel{\LeftR{\lforall}}
\UnaryInf$ \lforall[x][!A(x)],\Gamma \fCenter \Delta$
\DisplayProof
\hfill
\Axiom$ \Gamma \fCenter \Delta, !A(a) $
\RightLabel{\RightR{\lforall}}
\UnaryInf$ \Gamma \fCenter \Delta, \lforall[x][!A(x)]$
\DisplayProof
\end{defish}

في~\LeftR{\lforall} يكون~$t$ حدًا مغلقًا (أي حدًا بلا متغيرات). وفي
\RightR{\lforall} يكون~$a$ !!a{constant} لا يجوز أن يرد في أي موضع
من التتابعية السفلى لقاعدة~\RightR{\lforall}. ونسمي~$a$
\emph{المتغير المميَّز} لاستدلال~\RightR{\forall}.\footnote{نستعمل
مصطلح «المتغير المميَّز» مع أن~$a$ في القاعدة أعلاه !!a{constant}؛
ولذلك أسباب تاريخية.}

\subsection{قواعد $\lexists$}

\begin{defish}
\Axiom$ !A(a), \Gamma \fCenter \Delta $
\RightLabel{\LeftR{\lexists}}
\UnaryInf$ \lexists[x][!A(x)], \Gamma \fCenter \Delta$
\DisplayProof
\hfill
\Axiom$ \Gamma \fCenter \Delta, !A(t) $
\RightLabel{\RightR{\lexists}}
\UnaryInf$ \Gamma \fCenter \Delta, \lexists[x][!A(x)]$
\DisplayProof
\end{defish}

مرة أخرى، يكون~$t$ حدًا مغلقًا، ويكون~$a$ !!a{constant} لا يرد في
التتابعية السفلى لقاعدة~\LeftR{\lexists}. ونسمي~$a$
\emph{المتغير المميَّز} لاستدلال~\LeftR{\lexists}.

يسمى الشرط القاضي بألا يرد المتغير المميَّز في التتابعية السفلى
لاستدلال~\RightR{\lforall} أو~\LeftR{\lexists}
\emph{شرط المتغير المميَّز}.

\begin{explain}
تذكّر الاصطلاح القائل إننا، حين تكون~$!A$ !!a{formula} يكون فيها
ال!!{variable}~$x$ حرًا، نشير إلى ذلك بكتابة~$!A(x)$. وفي السياق
نفسه تكون~$!A(t)$ اختصارًا لـ~$\Subst{!A}{t}{x}$. ولذلك كان يمكننا
أيضًا كتابة قاعدة~\RightR{\lexists} كما يأتي:
\begin{prooftree}
  \Axiom$\Gamma \fCenter \Delta, \Subst{!A}{t}{x}$
  \RightLabel{\RightR{\lexists}}
  \UnaryInf$\Gamma \fCenter \Delta, \lexists[x][!A]$
\end{prooftree}
لاحظ أن~$t$ قد يرد سلفًا في~$!A$؛ فقد تكون~$!A$، مثلًا،
$\Atom{\Obj P}{t,x}$. ومن ثم فإن استنتاج
$\Gamma \Sequent \Delta, \lexists[x][\Atom{\Obj P}{t,x}]$ من
~$\Gamma \Sequent \Delta, \Atom{\Obj P}{t,t}$ تطبيق صحيح
لـ~\RightR{\lexists}---إذ يجوز لك أن «تستبدل» ورودًا واحدًا أو أكثر
لـ~$t$ في المقدمة بال!!{variable} المربوط~$x$، ولا يلزم استبدالها
كلها. غير أن شرطي المتغير المميَّز في~\RightR{\lforall}
و~\LeftR{\lexists} يقتضيان ألا يرد ال!!{constant}~$a$ في~$!A$.
ولذلك لا يصح أن تستنتج
$\Gamma \Sequent \Delta, \lforall[x][\Atom{\Obj P}{a,x}]$ من
$\Gamma \Sequent \Delta, \Atom{\Obj P}{a,a}$ باستعمال
~$\RightR{\lforall}$.
\end{explain}

\begin{explain}
لا ترد أي قيود على الحد~$t$ في~\RightR{\lexists}
و~\LeftR{\lforall}. أما في قاعدتي~\LeftR{\lexists}
و~\RightR{\lforall}، فيقتضي شرط المتغير المميَّز ألا يرد
ال!!{constant}~$a$ في أي موضع خارج~$!A(a)$ في التتابعية العليا.
وهذا ضروري لضمان سلامة النظام، أي إنه لا !!{derive}s إلا تتابعيات
صالحة. ومن دون هذا الشرط كان ما يأتي جائزًا:
\begin{prooftree}
  \Axiom$!A(a) \fCenter !A(a)$
  \RightLabel{*\LeftR{\lexists}}
  \UnaryInf$\lexists[x][!A(x)] \fCenter !A(a)$
  \RightLabel{\RightR{\lforall}}
  \UnaryInf$\lexists[x][!A(x)] \fCenter \lforall[x][!A(x)]$
  \DisplayProof\bottomAlignProof
  \qquad
  \Axiom$!A(a) \fCenter !A(a)$
  \RightLabel{*\RightR{\lforall}}
  \UnaryInf$!A(a) \fCenter \lforall[x][!A(x)]$
  \RightLabel{\LeftR{\lexists}}
  \UnaryInf$\lexists[x][!A(x)] \fCenter \lforall[x][!A(x)]$
\end{prooftree}
لكن التتابعية~$\lexists[x][!A(x)] \Sequent \lforall[x][!A(x)]$ ليست
صالحة.
\end{explain}

\end{document}
